% ==============================================================
% Section 5: Economic Interpretations
% ==============================================================

\section{Economic Interpretations}
\label{sec:interpretations}

The mathematical structure developed in the preceding sections---a CES objective, a linear constraint, an ideal price index, and the duality between value and cost---governs a family of economic problems. This section presents the taxonomy. Later chapters develop each application in detail; here we establish the common structure and vocabulary.

% --------------------------------------------------------------
\subsection{The Unified Framework}
\label{subsec:unified-framework}
% --------------------------------------------------------------

Every problem we have solved takes the form:
\begin{equation}
\max_{\bx} \; \M(\bx; \bom, \rho) \quad \text{subject to} \quad \sum_{i=1}^{n} p_i x_i = W,
\end{equation}
with solution characterized by:
\begin{align}
x_i\opt &= \omega_i \left( \frac{p_i}{\Pstar} \right)^{-\varepsilon} \cdot \frac{W}{\Pstar}, \\[4pt]
\Pstar &= \left( \sum_{j=1}^{n} \omega_j \, p_j^{1-\varepsilon} \right)^{\frac{1}{1-\varepsilon}}, \\[4pt]
\M\opt &= \frac{W}{\Pstar}.
\end{align}

The dual problem---minimizing $\sum_i p_i x_i$ subject to $\M(\bx) \geq \bar{m}$---yields the cost function $C = \Pstar \cdot \bar{m}$. The same price index $\Pstar$ mediates both problems.

What changes across applications is the \emph{interpretation} of symbols:

\begin{table}[H]
\centering
\caption{Taxonomy of CES Optimization Problems}
\label{tab:taxonomy}
\renewcommand{\arraystretch}{1.4}
\begin{tabular}{@{}l c c c c c@{}}
\toprule
\textbf{Application} & $x_i$ & $p_i$ & $W$ & $\omega_i$ & $\Pstar$ \\
\midrule
Consumer choice & consumption & goods prices & wealth & preferences & consumer price index \\[2pt]
Production & inputs & factor prices & cost budget & technology & unit cost \\[2pt]
Portfolio & state consumption & state prices & wealth & probabilities & risk-adjusted price \\[2pt]
Transformation & outputs & output prices & capacity & technology & output price index \\
\bottomrule
\end{tabular}
\end{table}

\begin{calloutimportant}[One Mathematics, Many Stories]
The optimal allocation formula, the price index, the duality relations, and the welfare results are \emph{identical} across all applications. Only the economic narrative differs. This is the power of the generalized mean framework: solve once, interpret many times.
\end{calloutimportant}

% --------------------------------------------------------------
\subsection{Consumer Choice}
\label{subsec:consumer-choice}
% --------------------------------------------------------------

The consumer maximizes utility over a bundle of goods:
\begin{equation}
\max_{\{c_i\}} \; \U(\bc; \bom, \rho) \quad \text{subject to} \quad \sum_{i=1}^{n} p_i c_i = W,
\end{equation}
where $c_i$ is consumption of good $i$, $p_i$ is its price, and $W$ is wealth.

\paragraph{Interpretation of Parameters.}
The weights $\omega_i$ reflect preference intensity---how much the consumer values good $i$ relative to others. The elasticity $\varepsilon = 1/(1-\rho)$ governs substitutability: how readily the consumer switches between goods as relative prices change.

\paragraph{The Consumer Price Index.}
The ideal price index $\Pstar$ is the \emph{cost-of-living index}---the minimum expenditure required to achieve one unit of utility. Indirect utility is $V = W / \Pstar$: real wealth measured in utility units.

\paragraph{Budget Shares.}
The optimal expenditure share on good $i$ is $w_i = p_i c_i / W$. These are the ``revealed'' weights that emerge from optimization, generally differing from the primitive weights $\omega_i$ when prices are heterogeneous.

\begin{calloutnote}[Demand Estimation]
In applied work, the CES demand system is estimated from expenditure data. The budget share equation
\[
w_i = \frac{\omega_i \, p_i^{1-\varepsilon}}{\sum_j \omega_j \, p_j^{1-\varepsilon}}
\]
is log-linear in prices after transformation, facilitating econometric estimation of $\varepsilon$ and the $\omega_i$. We return to estimation in a later chapter.
\end{calloutnote}

% --------------------------------------------------------------
\subsection{Production and Cost Minimization}
\label{subsec:production}
% --------------------------------------------------------------

A firm combines inputs to produce output using CES technology:
\begin{equation}
Y = \F(\bx; \bom, \phi) = \left( \sum_{i=1}^{n} \omega_i^{1-\phi} x_i^{\phi} \right)^{1/\phi},
\end{equation}
where $x_i$ is the quantity of input $i$ (labor, capital, materials), $\phi < 1$ is the technology parameter, and $\sigma = 1/(1-\phi)$ is the elasticity of substitution across inputs.

\paragraph{Cost Minimization.}
The firm minimizes cost $\sum_i r_i x_i$ subject to producing at least $\bar{Y}$, where $r_i$ is the price of input $i$. This is precisely the dual problem. The solution:
\begin{equation}
C(r, \bar{Y}) = \Pstar_r \cdot \bar{Y},
\end{equation}
where $\Pstar_r = (\sum_j \omega_j \, r_j^{1-\sigma})^{1/(1-\sigma)}$ is the \emph{unit cost}---the minimum cost to produce one unit of output.

\paragraph{Input Demands.}
Conditional input demands are:
\begin{equation}
x_i^c = \omega_i \left( \frac{r_i}{\Pstar_r} \right)^{-\sigma} \cdot \bar{Y}.
\end{equation}
These are ``compensated'' in the sense of holding output fixed, analogous to Hicksian demands in consumer theory.

\paragraph{Profit Maximization.}
If the firm is a price-taker in the output market at price $P$, it maximizes profit $P \cdot Y - \sum_i r_i x_i$. Substituting the cost function, profit is $(P - \Pstar_r) \cdot Y$ when $P > \Pstar_r$. The firm produces if and only if the output price exceeds unit cost.

\begin{calloutnote}[Returns to Scale]
The CES production function as written exhibits constant returns to scale: $\F(\lambda \bx) = \lambda \F(\bx)$. Variable returns can be introduced by raising $\F$ to a power, but we defer this to the chapter on nested structures.
\end{calloutnote}

% --------------------------------------------------------------
\subsection{Portfolio Choice and Risk}
\label{subsec:portfolio}
% --------------------------------------------------------------

Consider an investor allocating wealth across $n$ states of the world. State $i$ occurs with probability $\pi_i$, and the investor chooses consumption $c_i$ in each state.

\paragraph{Expected Utility.}
With CRRA (constant relative risk aversion) preferences, expected utility is:
\begin{equation}
\E[u(c)] = \sum_{i=1}^{n} \pi_i \frac{c_i^{1-\gamma}}{1-\gamma},
\end{equation}
where $\gamma > 0$ is the coefficient of relative risk aversion.

\paragraph{Certainty Equivalent.}
The certainty equivalent---the sure consumption that yields the same utility as the lottery---is:
\begin{equation}
\CE = \left( \sum_{i=1}^{n} \pi_i \, c_i^{1-\gamma} \right)^{\frac{1}{1-\gamma}} = \Mbar(\bc; \mathbf{p}i, 1-\gamma).
\end{equation}
This is a \emph{classical} generalized mean with power $\rho = 1 - \gamma$ and weights equal to probabilities.

\paragraph{Arrow-Debreu Prices.}
Let $q_i$ be the price of one unit of consumption in state $i$ (an Arrow-Debreu security). The budget constraint is $\sum_i q_i c_i = W$. Maximizing the certainty equivalent subject to this constraint yields:
\begin{equation}
c_i\opt = \pi_i^{1/\gamma} \left( \frac{q_i}{Q^*} \right)^{-1/\gamma} \cdot \frac{W}{Q^*},
\end{equation}
where $Q^* = (\sum_j \pi_j^{1/\gamma} q_j^{1-1/\gamma})^{\gamma/(\gamma-1)}$ is the risk-adjusted price index.

\begin{calloutimportant}[Risk-Neutral Probabilities]
The transformation from physical probabilities $\pi_i$ to the weights appearing in demands involves the same weight transformation $T_\rho$ from the chapter on generalized means. The resulting ``risk-neutral'' probabilities are central to asset pricing. We develop this connection fully in the chapter on uncertainty.
\end{calloutimportant}

% --------------------------------------------------------------
\subsection{Transformation and the PPF}
\label{subsec:transformation}
% --------------------------------------------------------------

Consider an economy that can produce multiple outputs subject to a transformation constraint:
\begin{equation}
\Hfun(y_1, \ldots, y_n; \boldsymbol{\eta}, h) = \left( \sum_{i=1}^{n} \eta_i^{1-h} y_i^{h} \right)^{1/h} \leq 1,
\end{equation}
where $y_i$ is output of good $i$, $\eta_i$ reflects productivity, and $h > 1$ makes the constraint convex (the feasible set is convex).

\paragraph{Revenue Maximization.}
A planner or firm maximizes revenue $\sum_i p_i y_i$ subject to the transformation constraint. As shown in Section~3.5, the solution involves the output price index:
\begin{equation}
Q^* = \left( \sum_{j=1}^{n} \eta_j \, p_j^{\frac{h}{h-1}} \right)^{\frac{h-1}{h}},
\end{equation}
and optimal outputs:
\begin{equation}
y_i\opt = \eta_i \left( \frac{p_i}{Q^*} \right)^{\frac{1}{h-1}}.
\end{equation}
Maximum revenue equals $R^* = Q^*$.

\paragraph{The Production Possibilities Frontier.}
The constraint $\Hfun(\by) = 1$ defines the PPF---the boundary of feasible output combinations. The curvature parameter $h$ governs the trade-off: higher $h$ means more convexity, reflecting increasing opportunity costs of transformation.

\paragraph{Connection to Consumer Problems.}
The revenue maximization problem is the ``producer dual'' to the consumer's expenditure minimization. Both involve a CES aggregator and a linear objective, but with roles reversed:

\begin{center}
\renewcommand{\arraystretch}{1.3}
\begin{tabular}{@{}lcc@{}}
\toprule
& \textbf{Consumer} & \textbf{Producer (PPF)} \\
\midrule
CES aggregator & Objective (concave, $\rho < 1$) & Constraint (convex, $h > 1$) \\
Linear function & Constraint (budget) & Objective (revenue) \\
Price index role & Deflates budget to utility & Inflates technology to revenue \\
\bottomrule
\end{tabular}
\end{center}

% --------------------------------------------------------------
\subsection{Summary: The Four Contexts}
\label{subsec:four-contexts}
% --------------------------------------------------------------

\Cref{tab:four-contexts} summarizes the key objects across all four applications.

\begin{table}[H]
\centering
\caption{Key Objects Across Applications}
\label{tab:four-contexts}
\renewcommand{\arraystretch}{1.5}
\small
\begin{tabular}{@{}l c c c c@{}}
\toprule
\textbf{Object} & \textbf{Consumer} & \textbf{Production} & \textbf{Portfolio} & \textbf{PPF} \\
\midrule
Quantities & $c_i$ (consumption) & $x_i$ (inputs) & $c_i$ (state cons.) & $y_i$ (outputs) \\
Prices & $p_i$ (goods) & $r_i$ (factors) & $q_i$ (state prices) & $p_i$ (outputs) \\
Wealth/Capacity & $W$ (wealth) & $C$ (cost) & $W$ (wealth) & $1$ (normalized) \\
Weights & $\omega_i$ (preferences) & $\omega_i$ (technology) & $\pi_i$ (probabilities) & $\eta_i$ (technology) \\
Price index & $\Pstar$ (CPI) & $\Pstar_r$ (unit cost) & $Q^*$ (risk-adj.) & $Q^*$ (output index) \\
Value/Cost & $V = W/\Pstar$ & $C = \Pstar_r \bar{Y}$ & $\CE = W/Q^*$ & $R^* = Q^*$ \\
Curvature & $\rho < 1$ (concave) & $\phi < 1$ (concave) & $1-\gamma < 1$ & $h > 1$ (convex) \\
\bottomrule
\end{tabular}
\end{table}

\begin{callouttip}[Reading the Table]
Each column is a complete economic model. Read vertically to see how abstract symbols map to concrete economic objects. Read horizontally to see how the same mathematical object (e.g., the price index) takes different economic meanings across contexts.
\end{callouttip}

% --------------------------------------------------------------
\subsection{Looking Ahead}
\label{subsec:looking-ahead}
% --------------------------------------------------------------

This taxonomy establishes the vocabulary; subsequent chapters develop each application in depth:

\paragraph{Consumer Theory.}
We examine demand elasticities, welfare measures (compensating and equivalent variation), aggregation across consumers, and the estimation of preference parameters from expenditure data.

\paragraph{Production and Trade.}
The CES production function underlies models of factor substitution, technical change, and international trade. We study multi-sector models, input-output linkages, and the gains from trade.

\paragraph{Risk and Uncertainty.}
The connection between CES and CRRA preferences opens the door to asset pricing, portfolio optimization, and the equity premium puzzle. The transformation between physical and risk-neutral probabilities is central.

\paragraph{General Equilibrium.}
When consumers and producers interact through markets, the four problems connect. Consumer demands meet producer supplies; factor payments fund consumer budgets; transformation possibilities constrain aggregate feasibility. The price indices from each problem must be mutually consistent in equilibrium.

\begin{calloutimportant}[The Recursive Structure]
Many economic models feature \emph{nested} CES aggregators---for example, a utility function over broad categories, each of which aggregates finer goods. The recursive structure preserves tractability while allowing richer substitution patterns. This is the subject of the chapter on nesting and recursion.
\end{calloutimportant}
